% ============================================================
%  SECTIONS 7, 2, AND 1
%  Updated: wuerstchen → wuerstchen, all numbers from NB10+NB13.
% ============================================================

\begin{abstract}
We introduce \textbf{\datasetname}, a 70{,}000-image benchmark for
AI-generated image detection comprising 10{,}000 real photographs
paired with 60{,}000 synthetic images produced by six architecturally
diverse text-to-image diffusion models: Stable Diffusion~1.5, SDXL,
FLUX.1-schnell, Kandinsky~2.2, PixArt-$\Sigma$, and
W\"{u}rstchen.
Unlike prior datasets, \textbf{\datasetname} applies a single
\emph{canonical preprocessing pipeline} --- JPEG equalisation followed
by lossless PNG serialisation --- identically to both real and AI
images, eliminating the JPEG compression and resolution asymmetries
that allow detectors to exploit format shortcuts rather than synthesis
artefacts.
All synthetic images are generated from captions derived from the real
images themselves, ensuring that semantic content is matched across
classes.
Splits are assigned at the pair level so that no real image and its
AI partners appear in different folds.
We validate the dataset with a one-epoch ResNet-18 sniff test,
confirm an overall detectability of 0.860 across architecturally
diverse generators, and provide baselines for three standard detectors.
EfficientNet-B4 achieves an AUROC of 0.9807 overall, while a
leave-one-out cross-generator experiment reveals a sharp
UNet-vs-non-UNet generalisation gap: detectors generalise poorly
to held-out UNet generators (AUROC~$\approx 0.86$--$0.89$) but
near-perfectly to held-out non-UNet generators
(AUROC~$\approx 0.99$).
The dataset, preprocessing code, generation notebooks, and baseline
checkpoints are publicly available at
\url{https://huggingface.co/datasets/Shanmuk4622/ai-detection-dataset-v2}.
\end{abstract}


\section{Introduction}
\label{sec:intro}

The capability of text-to-image generative models to produce
photorealistic images from natural language descriptions has improved
dramatically over the past three years.
Models such as Stable Diffusion~\cite{rombach2022ldm},
SDXL~\cite{podell2024sdxl}, and FLUX have reduced the cost of
synthetic image creation to near zero, enabling benign applications
in art and design while simultaneously expanding the threat surface
for disinformation, identity fraud, and manipulated media.
The urgency of reliable detection has grown in proportion: a recent
comprehensive evaluation of open-source detectors found that mean
detection accuracy against 2024-era generators has fallen to
approximately 38\%, down from around 79\% against generators from
2020--2021~\cite{bai2026opendetect}, illustrating an accelerating
arms race in which synthesis advances consistently outpace detection.

Constructing a detection dataset that is genuinely useful for this
problem is harder than it appears.
Two structural biases have been identified in widely used benchmarks.
The first is a \emph{JPEG compression disparity}: real images sourced
from the web are typically stored as JPEG with quality factors around
96, while generated images are saved as lossless PNG.
Detectors trained on such data learn to distinguish real from AI
primarily by measuring DCT coefficient distributions --- a shortcut
with no connection to synthesis artefacts and no robustness to
post-processing~\cite{grommelt2024jpeg}.
The second is a \emph{content distribution mismatch}: when prompts are
chosen independently of the real images, detectors may exploit
systematic differences in scene composition, object frequency, or
caption vocabulary rather than visual synthesis characteristics.
Both biases inflate reported detection numbers without producing
detectors that generalise to real-world conditions.

This paper addresses both problems simultaneously through three
design choices.
First, every image --- real or AI --- is passed through an identical
canonical preprocessing function that equalises JPEG compression
history and strips all metadata, making format-based shortcuts
structurally impossible.
Second, all synthetic images are generated from BLIP-2
captions~\cite{li2023blip2} derived directly from the real images,
so the semantic content available to every generator is held constant.
Third, splits are assigned at the \emph{pair level}: a real image and
all six of its AI partners always appear in the same fold, preventing
prompt content from leaking across the split boundary.
Figure~\ref{fig:pairing} illustrates the pairing and preprocessing
concept.

The resulting dataset, \textbf{\datasetname}, contains 70,000 images
spanning six generators that cover the three principal backbone
families in modern text-to-image synthesis: UNet latent diffusion
models~\cite{rombach2022ldm,podell2024sdxl}, attention-based Diffusion
Transformers~\cite{chen2024pixart}, and hierarchical two-stage
architectures~\cite{razzhigaev2023kandinsky,pernias2023wuerstchen}.
This architectural breadth is a deliberate contribution: prior paired
datasets~\cite{jiang2024passive} cover at most four generators from
a single family, making it impossible to measure how well a detector
trained on one architectural paradigm generalises to another.

\paragraph{Contributions.}
\begin{enumerate}[leftmargin=1.2em, itemsep=1pt]
\item A \textbf{70{,}000-image paired dataset} in which real and AI
  images share identical preprocessing, identical content grounding
  through shared BLIP-2 captions, and deterministic pair-level splits.
\item A \textbf{leak-free canonical preprocessing pipeline}
  (Algorithm~\ref{alg:preprocess}) that eliminates JPEG compression
  and resolution asymmetries between the two classes.
\item \textbf{Baseline detector results} for
  ResNet-50~\cite{he2016resnet},
  EfficientNet-B4~\cite{tan2019efficientnet}, and
  ViT-B/16~\cite{dosovitskiy2021vit}, including per-generator breakdown
  and a seven-way leave-one-out cross-generator generalisation
  experiment.
\item \textbf{Empirical evidence of a UNet-vs-non-UNet
  generalisation gap}: detectors generalise poorly to held-out UNet
  generators (AUROC~$\approx 0.86$--$0.89$) but near-perfectly to
  held-out non-UNet generators (AUROC~$\approx 0.99$), with
  implications for how future detection benchmarks should be structured.
\end{enumerate}

\paragraph{Paper organisation.}
Section~\ref{sec:related} surveys existing detection datasets and
methods.
Section~\ref{sec:methodology} describes the full dataset construction
pipeline.
Section~\ref{sec:audit} presents the leak audit and sniff-test
validation.
Section~\ref{sec:experiments} reports baseline detector results.
Section~\ref{sec:discussion} discusses the findings and limitations.
Section~\ref{sec:ethics} states ethical considerations and
reproducibility information.
Section~\ref{sec:conclusion} concludes.


\section{Discussion}
\label{sec:discussion}

\subsection{Why ResNet-50 Underperforms}
\label{ssec:disc_resnet}

The most striking result in Table~\ref{tab:baselines} is the gap
between ResNet-50 (AUROC~$= 0.7032$) and EfficientNet-B4
(AUROC~$= 0.9807$) despite both being trained under identical
conditions.
The gap is not explained by parameter count alone: EfficientNet-B4
has approximately 19~million parameters versus ResNet-50's 25~million,
yet achieves substantially higher discriminative performance.

The more likely explanation lies in the features each architecture
extracts.
ResNet-50's residual blocks use $3 \times 3$ convolutions effective
at capturing local texture but limited at modelling long-range spatial
dependencies.
Synthesis artefacts from modern diffusion models --- particularly
subtle high-frequency patterns in flat regions left by the denoising
process --- are not strictly local; they manifest as spatially
correlated noise requiring receptive fields larger than a single
convolutional block to detect reliably.
EfficientNet-B4's compound scaling increases width, depth, and input
resolution together, giving it larger effective receptive fields and
more capacity to model spatially extended artefact patterns.

\subsection{W\"{u}rstchen and the Detectability Gradient}
\label{ssec:disc_wuerstchen}

The per-generator results reveal a nuanced picture for W\"{u}rstchen.
In the sniff test it falls in the healthy zone (0.843), the only
generator to do so --- yet trained EfficientNet-B4 achieves AUROC
$= 0.9745$, confirming that its artefacts are detectable under
sufficient training.
This combination is unusual: a generator that is hard to detect
in a single-epoch probe but near-trivially detectable when held out
(cross-gen AUROC~$= 0.992$).

The explanation is that Würstchen's artefacts are distinct from all
UNet generators but not similar to real photographs in the feature
space that a multi-epoch trained detector learns.
In cross-generator generalisation, a detector that has never seen
W\"{u}rstchen is essentially told ``this image is unlike anything
I've been trained on'' --- and because all trained-on generators
are UNet-based or Transformer-based, a hierarchical Würstchen image
reads as anomalous, yielding near-perfect out-of-distribution
detection without direct exposure.

\subsection{The UNet-vs-Non-UNet Generalisation Gap}
\label{ssec:disc_crossgen}

The cross-generator generalisation experiment
(Section~\ref{ssec:cross_gen}) reveals that generalisation difficulty
depends almost entirely on backbone architecture.
UNet-based generators (SD~1.5, AUROC~$= 0.890$; SDXL,
AUROC~$= 0.861$) are substantially harder to detect when held out
than non-UNet generators (W\"{u}rstchen, AUROC~$= 0.992$;
Kandinsky~2.2, AUROC~$= 0.993$; PixArt-$\Sigma$, AUROC~$= 0.993$).

One mechanism is \emph{artefact distinctiveness}: non-UNet
architectures produce visual outputs with unusual frequency profiles
--- noted empirically by Corvi~\etal~\cite{corvi2023detection} ---
sufficiently different from any UNet output that a detector identifies
them as out-of-distribution even without direct exposure.
UNet generators produce artefacts similar across the family, so a
detector trained on some UNet generators learns a representation that
partially describes a held-out UNet, but not distinctively enough for
high precision (F1 scores of 0.667 and 0.657 for SD~1.5 and SDXL).

A dataset containing only UNet-based generators --- which describes
most existing benchmarks --- cannot reveal this gap and will
systematically overestimate generalisation when a non-UNet generator
is deployed.
\textbf{\datasetname}'s coverage of all three backbone families makes
this measurement possible for the first time in a paired, leak-free
setting.

\subsection{Limitations}
\label{ssec:disc_limits}

\paragraph{Scale.}
At 70,000 images, \textbf{\datasetname} is smaller than
GenImage~\cite{zhu2023genimage} (1.3~million) or
AI-GenBench~\cite{pellegrini2025aigenbench} (310,000).
The trade-off was explicit: preprocessing integrity and content
matching over raw scale.

\paragraph{Text-to-image scope.}
The dataset covers text-to-image generation only.
Detectors trained on \textbf{\datasetname} target the text-to-image
threat model and should not be assumed to generalise to editing-based
or reference-conditioning synthesis.

\paragraph{Inference hyperparameters.}
Step counts reflect real-world fast-inference usage.
Images at higher step counts may carry different artefact profiles.

\paragraph{Single captioner.}
All prompts were generated by BLIP-2
(\texttt{blip2-opt-2.7b})~\cite{li2023blip2}.
Richer or more diverse captions would increase prompt diversity.

\paragraph{ViT per-generator breakdown.}
Per-generator AUROC for ViT-B/16 was not computed due to memory
constraints; only overall metrics are reported.


\section{Conclusion}
\label{sec:conclusion}

We presented \textbf{\datasetname}, a 70,000-image paired dataset for
AI-generated image detection built around two principles that most
existing benchmarks violate: a canonical preprocessing pipeline
applied identically to real and AI images, and caption-grounded
content matching that eliminates semantic distribution differences
between classes.
Across six architecturally diverse generators spanning UNet, Diffusion
Transformer, and hierarchical backbone families, EfficientNet-B4
achieves an AUROC of 0.9807 under bias-free conditions.

The cross-generator generalisation experiment revealed a consistent
architectural divide: detectors generalise poorly to unseen UNet-based
generators (AUROC~$\approx 0.86$--$0.89$) while generalising
near-perfectly to unseen non-UNet generators (AUROC~$\approx 0.99$).
This gap is invisible to single-family datasets and has direct
implications for how detection benchmarks should be structured as
new generator families emerge.

The complete dataset, eleven-notebook build pipeline, all baseline
checkpoints, and validation report are released at
\url{https://huggingface.co/datasets/Shanmuk4622/ai-detection-dataset-v2}.

%\bibitem{pernias2023wuerstchen}
%P.~Pernias, D.~Rampas, M.~L.~Richter, C.~Pal, and M.~Aubreville,
%``W\"{u}rstchen: An efficient architecture for large-scale
%text-to-image diffusion models,''
%in \emph{Proc. Int. Conf. Learning Representations (ICLR)}, 2024.

%\bibitem{bai2026opendetect}
%H.~Bai \etal,
%``How well are open-sourced AI-generated image detection models
%out-of-the-box: A comprehensive benchmark study,''
%\emph{arXiv preprint arXiv:2602.07814}, 2026.
